\documentclass[10pt,a4paper]{amsart}
\usepackage[margin=19mm]{geometry}
\usepackage{amsmath,amssymb,mathtools}
\usepackage[scr=rsfs,cal=euler]{mathalfa}
\usepackage{xfrac}
\usepackage[unicode,hidelinks,bookmarks=false]{hyperref}
\newcommand{\ie}{i.e.\ }
\newcommand{\cf}{cf.\ }
\newcommand{\resp}{respectively\ }
\newcommand{\Subsection}{\subsection}
\providecommand{\nobreakdash}{\nobreak\hbox{-}}
\setlength{\emergencystretch}{2em}
\multlinegap0pt
\usepackage{amssymb}
\usepackage{mathtools}
\usepackage{enumerate}
\usepackage{xcolor}
\usepackage{tcolorbox}
\newtheorem{assumption}{Assumption}
\newtheorem{lemma}{Lemma}
\newcommand{\E}{\mathbb{E}}
\renewcommand{\P}{\mathbb{P}}
\newcommand{\W}{\mathcal{W}}
\title{On the structure of marginals in high dimensions}
\author{Daniel Bartl, Shahar Mendelson}
\date{Source: arXiv:2603.17291v1 (CC BY 4.0)}
\begin{document}
\maketitle

\noindent\emph{Source and licence.} On the structure of marginals in high dimensions. Version arXiv:2603.17291v1 (CC BY 4.0), distributed by arXiv under the Creative Commons Attribution 4.0 International licence, http://creativecommons.org/licenses/by/4.0/, as recorded in the arXiv licence metadata for this exact version and shown on the abstract page https://arxiv.org/abs/2603.17291v1. Full licence notice: \texttt{licenses/arxiv-2603.17291-CC-BY-4.0.txt}.\par\smallskip

\noindent\textit{Editorial scope.} Counted unit: the article's lemma lem:reduction.to.truncation (source0.tex lines 555--562). The article's complete original proof of it is delivered with it (source0.tex lines 563--624, one \texttt{proof} environment of 62 lines). No mathematics is written, repaired or re-derived: the statement and the proof are the article's own lines, in the article's own notation, and no retained line is substituted or re-flowed.  Retained uncounted as the source-local context the proof needs: lines 230--234; lines 412--417; lines 481--486.  Nothing was omitted from any retained span, so the package carries no omission record.  No retained line contains a citation, so the wrapper carries no bibliography and no bibliography entry.  No retained line contains a figure environment, an image-inclusion command or drawing code, so no image is delivered and the package declares no asset dependency.  Editorial formatting only, in addition: an AMS \texttt{amsart} preamble that restates the article's own declarations and macros used by the retained lines, namely the environments \texttt{assumption}, \texttt{lemma} on separate counters, together with the standard packages those lines need.  The retained environments print in this document as Assumption 1, Lemma 1 in order of appearance, whereas the article prints the same items with the numbering of its own full document.  The complete native source of the article as distributed in the arXiv e-print for this exact version is bundled unchanged as \texttt{sources/196612/source0.tex}.
\begin{assumption}
\label{ass:main}
	There are  $\theta > \E \sup_{f\in\mathcal{F}} G_f $, $B\geq 1$ and an event $\Omega_{0}$ on which, for every $f,h\in\mathcal{F}\cup\{0\}$, 
	\[ \|f-h\|_{L_{2}(\P_{N})} \leq B   \|f-h\|_{L_{2}}  + \frac{\theta}{\sqrt N}.\]	
\end{assumption}

\begin{lemma}
\label{lem:W.f.plus.h}
	For every two measurable functions $f$ and $h$,
	\[ \W_{1}(F_{N,f+h}, F_{f+h} )
	\leq  \W_{1}(F_{N,f}, F_{f} ) +  \|h\|_{L_{1}(\P_{N})} + \|h\|_{L_{1}} . \] 
\end{lemma}

\begin{align}
\label{eq:almost.optioma.admissible.seq}
 \sup_{f\in\mathcal{F}} \sum_{s\geq 0} 2^{s/2} \| \Delta_{s} f \|_{L_{2}}
\leq 4 \gamma_{2}(\mathcal{F}^{\theta})
\leq 4 \gamma_{2}(\mathcal{F}).
\end{align}

\begin{lemma}
\label{lem:reduction.to.truncation}
There is an absolute constant $c$ such that for $(X_{i})_{i=1}^{N}\in\Omega_{0}$ and  every $f\in\mathcal{F}$,
\begin{align*}
 \W_{1}(F_{N,\pi_{s_{1}}f}, F_{\pi_{s_{1}}f}) 
\leq \W_{1}(F_{N,\mathcal{T}(f)}, F_{\mathcal{T}(f)} ) + cB^{2}\left(  d_{\mathcal{F}} \sqrt\Delta + \frac{\gamma_{2}(\mathcal{F})}{\sqrt N}\right).
\end{align*} 
\end{lemma}

\begin{proof}
By Lemma \ref{lem:W.f.plus.h},
\begin{align*}
 \W_{1}(F_{N,\pi_{s_{1}}f}, F_{\pi_{s_{1}}f}) 
&\leq \W_{1}(F_{N,\mathcal{T}(f)}, F_{\mathcal{T}(f)} ) \\
&\qquad + \left\| \pi_{s_1}f- \mathcal{T}(f) \right\|_{L_{1}(\P_N )} + \left\|\pi_{s_1} f-\mathcal{T}(f)\right\|_{L_{1}}.
\end{align*}

To estimate $\|\pi_{s_1}f-\mathcal{T}(f) \|_{L_{1}}$, set $b_{s}= \|\Delta_{s} f\|_{L_{2}}\sqrt{N/2^s}$ and by the  Cauchy-Schwarz inequality followed by  Markov's inequality,
\begin{align*}
 \left\| \psi_{s}(\Delta_{s}f) \right\|_{L_{1}}
&\leq\E\left[ |\Delta_{s} f | 1_{\{ |\Delta_{s} f |\geq b_{s} \} } \right] \\
&\leq \|\Delta_{s} f \|_{L_{2}} \sqrt{ \P( |\Delta_{s} f |\geq b_{s}) }  \\
&\leq \|\Delta_{s} f \|_{L_{2}} \sqrt{ \frac{2^{s}}{N} }.
\end{align*}
In a similar manner, using that $2^{s_{0}}\geq  \Delta N$ and that $\|\pi_{s_{0}} f\|_{L_{2}}\leq d_{\mathcal{F}} $,
\begin{align*}
 \left\| \psi_{s_{0}}(\pi_{s_{0}}f) \right\|_{L_{1}}
&\leq \| \pi_{s_{0}} f \|_{L_{2}} \sqrt{ \P( |\pi_{s_{0}} f |\geq \|\pi_{s_{0}} f\|_{L_{2}}/\sqrt{\Delta}) }  
\leq   d_{\mathcal{F}}  \sqrt\Delta.
\end{align*}
Thus, by \eqref{eq:almost.optioma.admissible.seq},
\begin{align}
\label{eq:pi.s1.vs.T.some.eq}
\begin{split}
\| \pi_{s_1} f-\mathcal{T}(f) \|_{L_{1}}
&\leq \left\| \psi_{s_{0}}(\pi_{s_{0}}f) \right\|_{L_{1}} + \sum_{s=s_{0}}^{s_{1}-1} \left\| \psi_{s}(\Delta_{s}f) \right\|_{L_{1}}\\
&\leq  d_{\mathcal{F}} \sqrt\Delta  + \sum_{s=s_{0}}^{s_{1}-1}  \sqrt\frac{2^{s}}{N} \left\|\Delta_{s}f \right\|_{L_{2}}
\leq d_{\mathcal{F}}  \sqrt\Delta  + 4 \frac{\gamma_{2}(\mathcal{F})}{\sqrt N}.
\end{split}
\end{align}
%where the second inequality follows because $(\mathcal{F}_{s}^{\theta})_{s\geq 0}$ is an almost optimal admissible sequence, see \eqref{eq:almost.optioma.admissible.seq}.


The analysis of the  $\|\psi_{s}(\Delta_{s} f)\|_{L_{1}(\P_N)}$-term is slightly more involved.
Since $\mathcal{F}_{s_{1}}^{\theta}$ is a $\theta/\sqrt{N}$-separated set, we either have $\Delta_{s} f=0$ (in which case  $\| \psi_{s}(\Delta_sf)\|_{L_{1}(\P_{N})} =0$), or $\|\Delta_{s } f\|_{L_{2}}\geq \theta/\sqrt N$.
In the latter case, and since $|\psi_{s}(\Delta_sf)| = \max\{0,|\Delta_sf|-b_{s}\}$,  tail-integration followed by Markov's inequality show that 
\begin{align*}
 \left\| \psi_{s}(\Delta_sf) \right\|_{L_{1}(\P_{N})}
& =\int_{b_{s}}^{\infty}\P_{N}(|\Delta_{s} f|> t) \,dt 
\leq \int_{b_{s}}^{\infty} \frac{ \|\Delta_{s} f \|_{L_2(\P_N)}^{2}}{ t^{2}} \,dt .
\end{align*}
Moreover, $\| \Delta_s f\|_{L_{2}}\geq\theta/\sqrt N$ and it thus follows from Assumption \ref{ass:main} that 
\[ \|\Delta_{s} f \|_{L_2(\P_N)}^{2} 
\leq  (B+1)^{2}  \|\Delta_{s} f \|_{L_2}^{2} .\]
Therefore, by the choice of $b_{s}$,
\begin{align*}
\left\| \psi_{s}(\Delta_sf) \right\|_{L_{1}(\P_{N})}
&\leq (B+1)^{2}  \|\Delta_{s} f \|_{L_2}^{2}  \int_{b_{s}}^{\infty} \frac{1}{t^{2}} \,dt \\
&=(B+1)^{2} \|\Delta_{s} f \|_{L_2} \sqrt{\frac{2^{s}}{N}}.
\end{align*}

The same arguments can be used to show that
\begin{align*}
 \left\| \psi_{s_{0}}(\pi_{s_{0}}f) \right\|_{L_{1}(\P_{N})}
&\leq c_{1}B^{2}  d_{\mathcal{F}}  \sqrt\Delta,
\end{align*}
and following the same path as in \eqref{eq:pi.s1.vs.T.some.eq}, 
\[\| \pi_{s_1}f-\mathcal{T}(f) \|_{L_{1}(\P_{N})} 
\leq c_{2} B^{2} \left(  d_{\mathcal{F}}  \sqrt\Delta  +  \frac{\gamma_{2}(\mathcal{F})}{\sqrt N} \right),\]
as required. 
\end{proof}

\end{document}
